% ============================================================================
% LANGENTWURF LE-02b: tener + edad (Alter angeben)
% Spanisch Klasse 7 - Sequenz 2: Mi familia
% ============================================================================
% Tier: 1 (vollständig)
% Doppelstunde: 2b (Std. 3-4 von 8)
% Fachfarbe: #c41e3a (Karminrot)
% ============================================================================

\documentclass[11pt,a4paper]{article}

\usepackage[utf8]{inputenc}
\usepackage[T1]{fontenc}
\usepackage[ngerman]{babel}
\usepackage{geometry}
\usepackage{fancyhdr}
\usepackage{lastpage}
\usepackage{xcolor}
\usepackage{tcolorbox}
\usepackage{enumitem}
\usepackage{tabularx}
\usepackage{longtable}
\usepackage{array}
\usepackage{booktabs}
\usepackage{amssymb}

% ============================================================================
% KONFIGURATION
% ============================================================================

\newcommand{\fach}{Spanisch}
\newcommand{\fachfarbe}{c41e3a}
\newcommand{\jahrgangsstufe}{7}
\newcommand{\schulart}{Gesamtschule}
\newcommand{\stundenthema}{tener + edad (Alter angeben)}
\newcommand{\reihenthema}{Mi familia}
\newcommand{\stundennummer}{2b}
\newcommand{\reihenstunden}{4}
\newcommand{\gerniveau}{A1}
\newcommand{\lehrwerk}{Qué pasa 1}
\newcommand{\lehrwerkseite}{S. 24-29}
\newcommand{\lerngruppengroesse}{24}

% ============================================================================
% LAYOUT
% ============================================================================
\geometry{left=2.5cm, right=2.5cm, top=2.5cm, bottom=2.5cm}
\definecolor{fach}{HTML}{\fachfarbe}
\definecolor{gray}{HTML}{666666}
\definecolor{lightgray}{HTML}{f5f5f5}
\definecolor{successgreen}{HTML}{2a7a4b}
\definecolor{warningorange}{HTML}{b35c00}
\definecolor{basis}{HTML}{2a7a4b}
\definecolor{standard}{HTML}{2c5aa0}
\definecolor{erweiterung}{HTML}{7b1fa2}

\tcbuselibrary{skins,breakable}

\newtcolorbox{infobox}[1][]{
    colback=lightgray,
    colframe=fach,
    fonttitle=\bfseries,
    title=#1,
    breakable
}

\newtcolorbox{scorebox}{
    colback=white,
    colframe=fach,
    fonttitle=\bfseries,
    title=Didaktik-Score,
    breakable
}

\pagestyle{fancy}
\fancyhf{}
\fancyhead[L]{\small\textcolor{gray}{\fach{} -- Jg. \jahrgangsstufe{} -- \stundenthema}}
\fancyhead[R]{\small\textcolor{gray}{LE-\stundennummer{} | Tier 1}}
\fancyfoot[C]{\small\thepage{} / \pageref{LastPage}}
\renewcommand{\headrulewidth}{0.4pt}

% Differenzierungsmarker
\newcommand{\B}{\textcolor{basis}{\textbf{[B]}}}
\newcommand{\Std}{\textcolor{standard}{\textbf{[S]}}}
\newcommand{\E}{\textcolor{erweiterung}{\textbf{[E]}}}

% ============================================================================
% DOKUMENT
% ============================================================================
\begin{document}

% ============================================================================
% DECKBLATT
% ============================================================================
\begin{titlepage}
    \centering
    \vspace*{2cm}
    {\large\textcolor{gray}{\schulart}}\\[2cm]
    {\Huge\textcolor{fach}{\textbf{Langentwurf}}}\\[0.5cm]
    {\large LE-\stundennummer{} | Tier 1 (vollständig)}\\[2cm]
    {\LARGE\textbf{\stundenthema}}\\[0.5cm]
    {\large Sequenz: \reihenthema}\\[2cm]
    \begin{tabular}{rl}
        \textbf{Fach:} & \fach{} (\gerniveau) \\[0.3cm]
        \textbf{Klasse:} & \jahrgangsstufe \\[0.3cm]
        \textbf{Lehrwerk:} & \lehrwerk, \lehrwerkseite \\[0.3cm]
        \textbf{Doppelstunde:} & \stundennummer{} (Std. 3-4 von 8) \\[0.3cm]
    \end{tabular}
    \vfill
\end{titlepage}

\tableofcontents
\newpage

% ============================================================================
% 1. BEDINGUNGSANALYSE
% ============================================================================
\section{Bedingungsanalyse}

\subsection{Lerngruppe}
\begin{tabular}{ll}
    Kursgröße: & \lerngruppengroesse{} SuS \\
    Jahrgangsstufe: & \jahrgangsstufe \\
    Sprachniveau: & \gerniveau{} (GER) -- Anfänger \\
    Fremdsprache: & 2. Fremdsprache (nach Englisch) \\
\end{tabular}

\subsection{Vorwissen aus 02a}
\begin{itemize}
    \item Familienmitglieder: padre, madre, hermano, hermana, abuelo, abuela
    \item Possessivbegleiter: mi, tu, su
    \item Verb \textit{ser} (aus Sequenz 1)
    \item Satzstruktur: \textit{Mi hermano se llama...}
\end{itemize}

\subsection{Differenzierung (intern)}
\begin{tabular}{|l|p{3.5cm}|p{3.5cm}|p{3.5cm}|}
\hline
& \B{} \textbf{Basis} & \Std{} \textbf{Standard} & \E{} \textbf{Erweiterung} \\
\hline
\textbf{Ziel} & BR: Rezeption & MR: Produktion mit Hilfe & GY: Freie Produktion \\
\hline
\textbf{tener} & Lückentext mit Vorgabe & Konjugation selbst & + Sätze bilden \\
\hline
\textbf{Zahlen} & Zuordnen (Wort-Zahl) & Zahlen schreiben & + Rechenaufgaben \\
\hline
\textbf{Alter} & Formel anwenden & Familie beschreiben & Stammbaum erstellen \\
\hline
\end{tabular}

\vspace{0.5em}
\textbf{WICHTIG:} Diese Marker sind NUR für die Lehrkraft!

% ============================================================================
% 2. SACHANALYSE
% ============================================================================
\section{Sachanalyse}

\subsection{Sprachliche Strukturen}

\textbf{Verb \textit{tener} (Präsens):}
\begin{center}
\begin{tabular}{|l|l|l|}
\hline
\textbf{Person} & \textbf{Pronomen} & \textbf{tener} \\
\hline
1. Sg. & yo & tengo \\
2. Sg. & tú & tienes \\
3. Sg. & él/ella/usted & tiene \\
1. Pl. & nosotros/-as & tenemos \\
2. Pl. & vosotros/-as & tenéis \\
3. Pl. & ellos/ellas/ustedes & tienen \\
\hline
\end{tabular}
\end{center}

\textbf{Zahlen 1-20:}
\begin{center}
\begin{tabular}{|l|l|l|l|l|}
\hline
1 uno & 5 cinco & 9 nueve & 13 trece & 17 diecisiete \\
2 dos & 6 seis & 10 diez & 14 catorce & 18 dieciocho \\
3 tres & 7 siete & 11 once & 15 quince & 19 diecinueve \\
4 cuatro & 8 ocho & 12 doce & 16 dieciséis & 20 veinte \\
\hline
\end{tabular}
\end{center}

\textbf{Redemittel -- Alter:}
\begin{itemize}
    \item \textbf{Frage:} ¿Cuántos años tienes? / ¿Cuántos años tiene [tu hermano]?
    \item \textbf{Antwort:} Tengo [número] años. / Mi hermano tiene [número] años.
\end{itemize}

\subsection{Kontrastive Betrachtung}

\begin{tabular}{|l|p{5cm}|p{5cm}|}
\hline
& \textbf{Deutsch} & \textbf{Spanisch} \\
\hline
Alter & \textit{sein} (Ich \textbf{bin} 12 Jahre alt) & \textit{tener} (Ich \textbf{habe} 12 Jahre) \\
\hline
Unregelmäßigkeit & -- & tengo (1. Sg.) mit Stammvokalwechsel \\
\hline
Zahlen 16-19 & Einzelwörter & Komposita: dieci + sechs = dieciséis \\
\hline
\end{tabular}

\textbf{Typische Interferenz:} *``Yo soy doce años'' (dt. Struktur) statt ``Tengo doce años''

\subsection{Kulturelle Aspekte}
\begin{itemize}
    \item Geburtstage in spanischsprachigen Ländern: Piñata, Quinceañera (15. Geburtstag)
    \item Höflichkeitsfrage nach dem Alter: In Spanien weniger tabu als in Deutschland
\end{itemize}

% ============================================================================
% 3. DIDAKTISCHE ANALYSE
% ============================================================================
\section{Didaktische Analyse}

\subsection{Stellung in der Sequenz}
\begin{tabular}{|c|l|l|}
\hline
\textbf{Block} & \textbf{Thema} & \textbf{Status} \\
\hline
2a & Familienmitglieder + Possessivbegleiter & abgeschlossen \\
\textbf{2b} & \textbf{tener + edad (Alter angeben)} & \textbf{diese Stunde} \\
2c & Adjektive + Beschreibung & folgt \\
2d & Mi familia (Präsentation) + Test & folgt \\
\hline
\end{tabular}

\subsection{Kompetenzziele (Rahmenplan MV)}
\begin{itemize}
    \item \textbf{Kommunikativ:} Sprechen (Alter erfragen und angeben), Schreiben (Familienbeschreibung)
    \item \textbf{Sprachlich:} Grammatik (Verb tener), Wortschatz (Zahlen 1-20)
    \item \textbf{Interkulturell:} Geburtstagsrituale in hispanischen Kulturen
    \item \textbf{Methodisch:} Musterdialoge als Gerüst nutzen
\end{itemize}

\subsection{Legitimation}
\textbf{Rahmenplan Spanisch MV (Kl. 7-10):}
\begin{itemize}
    \item Kompetenzbereich: Kommunikative Kompetenz -- Sprechen, Schreiben
    \item Standard: ``Die SuS können über sich und ihre Familie berichten und einfache Informationen austauschen.''
\end{itemize}

% ============================================================================
% 4. STUNDENZIELE
% ============================================================================
\section{Stundenziele}

\subsection{Grobziel}
Die SuS erweitern ihre \textbf{kommunikative Kompetenz im Bereich Sprechen und Schreiben}, indem sie das Verb \textit{tener} konjugieren, die Zahlen 1-20 anwenden und das Alter von Personen erfragen und angeben.

\subsection{Feinziele (operationalisiert)}
\begin{tabular}{|c|p{7.5cm}|c|c|}
\hline
\textbf{Nr.} & \textbf{Die SuS können...} & \textbf{AFB} & \textbf{Diff.} \\
\hline
FZ1 & die Formen von \textit{tener} (tengo, tienes, tiene, tenemos) \textbf{nennen} und korrekt aussprechen & I & alle \\
\hline
FZ2 & die Zahlen 1-20 auf Spanisch \textbf{benennen} und der Ziffer \textbf{zuordnen} & I & alle \\
\hline
FZ3 & das eigene Alter mit \textit{Tengo ... años} \textbf{angeben} und die Frage \textit{¿Cuántos años tienes?} \textbf{verstehen} und \textbf{beantworten} & II & alle \\
\hline
FZ4 & das Alter von Familienmitgliedern in Sätzen \textbf{beschreiben} (z.B. ``Mi hermana tiene 15 años'') & II & \Std{}/\E{} \\
\hline
FZ5 & einen kurzen Text über die Familie mit Altersangaben \textbf{verfassen} & III & \E{} \\
\hline
\end{tabular}

\vspace{1em}
\textbf{AFB-Verteilung:} AFB I: 25\% | AFB II: 50\% | AFB III: 25\%

% ============================================================================
% 5. METHODISCHE ANALYSE
% ============================================================================
\section{Methodische Analyse}

\subsection{Einstieg}
\begin{itemize}
    \item \textbf{Methode:} Geburtstagskalender der Klasse -- Wer hat wann Geburtstag?
    \item \textbf{Begründung:} Alltagsbezug, intrinsische Motivation (eigenes Alter)
    \item \textbf{Alternative:} Kurzvideo: Geburtstagslied ``Cumpleaños feliz''
\end{itemize}

\subsection{Erarbeitung}
\begin{itemize}
    \item \textbf{Methode:} Zahlenreihe 1-20 mit Bewegung (TPR: Total Physical Response)
    \item \textbf{Begründung:} Multisensorisches Lernen, Aktivierung
    \item \textbf{Scaffolding:} Zahlenplakat, Farbe für 11-15 (Sonderformen), Rhythmus
\end{itemize}

\subsection{Übungsphasen}
\begin{itemize}
    \item \textbf{Methode:} Zahlen-Bingo $\rightarrow$ Partnerarbeit: Alter erfragen $\rightarrow$ Familienbeschreibung
    \item \textbf{Progression:} Rezeptiv (Bingo) $\rightarrow$ Produktiv gelenkt (Dialog) $\rightarrow$ Produktiv frei (Text)
    \item \textbf{Differenzierung:} Hilfekarten mit Satzmustern für \B{}, Erweiterungsaufgaben für \E{}
\end{itemize}

\subsection{Medien}
\begin{itemize}
    \item Tafel/Whiteboard: tener-Konjugation, Zahlenreihe
    \item Lehrwerk: \lehrwerk, \lehrwerkseite
    \item Arbeitsblatt: AB-02b.pdf
    \item Lernpfad: LP-02b.html (Tablets, optional)
    \item Bingo-Karten (Zahlen 1-20)
\end{itemize}

% ============================================================================
% 6. VERLAUFSPLANUNG
% ============================================================================
\section{Verlaufsplanung}

\textbf{1. Stunde (45 Min.)}

\begin{longtable}{|p{1cm}|p{1.5cm}|p{4cm}|p{3cm}|p{1.5cm}|p{2cm}|}
\hline
\textbf{Zeit} & \textbf{Phase} & \textbf{Lehrerhandeln} & \textbf{Schülerhandeln} & \textbf{SF} & \textbf{Material} \\
\hline
\endhead

0--5 & Einstieg & Begrüßung (span.), Frage: ``¿Cuántos años tienes?'' mit Gestik & Vermutungen, erste Antworten & Plenum & -- \\
\hline

5--15 & Input & Zahlen 1-10 einführen mit Fingerzeigen, chorisches Sprechen & Nachsprechen, mitzählen & Plenum & Tafel \\
\hline

15--25 & Übung 1 & Zahlen 11-20 mit Rhythmus-Klatsch, Besonderheiten (11-15) hervorheben & Mitsprechen, klatschen & Plenum & Zahlenplakat \\
\hline

25--35 & Übung 2 & Zahlen-Bingo anleiten, Zahlen vorlesen & Bingo spielen, Zahlen erkennen & EA & Bingo-Karten \\
\hline

35--40 & Sicherung & Schnell-Quiz: ``¿Qué número es?'' & Antworten rufen & Plenum & Tafel \\
\hline

40--45 & Überleitung & Ausblick: tener + Alter; AB Merkkasten Zahlen ausfüllen & Notieren, AB beginnen & Plenum & AB-02b \\
\hline
\end{longtable}

\vspace{1em}

\textbf{2. Stunde (45 Min.)}

\begin{longtable}{|p{1cm}|p{1.5cm}|p{4cm}|p{3cm}|p{1.5cm}|p{2cm}|}
\hline
\textbf{Zeit} & \textbf{Phase} & \textbf{Lehrerhandeln} & \textbf{Schülerhandeln} & \textbf{SF} & \textbf{Material} \\
\hline
\endhead

0--5 & Aktivier. & Schnelles Zahlen-Quiz (Wiederholung) & Antworten & Plenum & -- \\
\hline

5--15 & Input & Verb \textit{tener} einführen: ``Yo tengo 12 años''; Konjugation an Tafel & Tabelle ergänzen (AB) & Plenum & Tafel, AB \\
\hline

15--25 & Übung 1 & PA: Alter erfragen ``¿Cuántos años tienes?'' -- ``Tengo ... años'' & Dialoge üben & PA & -- \\
\hline

25--35 & Übung 2 & AB Aufgabe 3 + 4: Alter angeben, Familie beschreiben & AB bearbeiten & EA/PA & AB-02b \\
\hline

35--40 & Sicherung & Ergebnisse besprechen, Fehlerkorrektur & Vergleichen, korrigieren & Plenum & Tafel \\
\hline

40--45 & Abschluss & Zusammenfassung, HA: Lernpfad LP-02b.html, Vokabeln lernen & Notieren & Plenum & -- \\
\hline
\end{longtable}

\textbf{Legende:} EA = Einzelarbeit, PA = Partnerarbeit

% ============================================================================
% 7. ERWARTETE SCHWIERIGKEITEN
% ============================================================================
\section{Erwartete Schwierigkeiten}

\begin{tabular}{|p{4.5cm}|p{8cm}|}
\hline
\textbf{Schwierigkeit} & \textbf{Reaktion} \\
\hline
Interferenz: *``Yo soy 12 años'' & Expliziter Hinweis: Spanisch = ``Jahre haben'' (tener), nicht ``Jahre sein'' \\
\hline
Verwechslung \textit{tiene/tienes} & Kontrastive Übung: ``Du hast'' (tienes) vs. ``Er/Sie hat'' (tiene) \\
\hline
Zahlen 11-15 (Sonderformen) & Farbliche Hervorhebung, Rhythmus-Lernen, Wiederholung \\
\hline
Aussprache \textit{dieciséis} & Silbenweise sprechen: die-ci-séis, chorisches Üben \\
\hline
Zeitproblem & Puffer: Bingo kürzen, Aufgabe 4 als HA \\
\hline
\end{tabular}

% ============================================================================
% 8. HAUSAUFGABE
% ============================================================================
\section{Hausaufgabe}

\begin{itemize}
    \item Lehrwerk S. 26, Übung 4 (tener einsetzen)
    \item Vokabeln lernen: Zahlen 1-20, tener-Formen
    \item Optional: Lernpfad LP-02b.html bearbeiten
    \item \E{}: Kurzen Text über die Familie schreiben (3-4 Sätze mit Altersangaben)
\end{itemize}

% ============================================================================
% 9. LÖSUNGEN
% ============================================================================
\section{Lösungen}

\subsection{AB-02b Lösungen}

\textbf{Aufgabe 1: tener einsetzen (4P)}
\begin{enumerate}[label=\alph*)]
    \item Yo \textbf{tengo} 12 años.
    \item Mi hermana \textbf{tiene} 15 años.
    \item ¿Cuántos años \textbf{tienes} tú?
    \item Nosotros \textbf{tenemos} un perro.
\end{enumerate}

\textbf{Aufgabe 2: Zahlen zuordnen (4P)}
\begin{itemize}
    \item doce $\rightarrow$ 12
    \item diecisiete $\rightarrow$ 17
    \item ocho $\rightarrow$ 8
    \item quince $\rightarrow$ 15
\end{itemize}

\textbf{Aufgabe 3: Alter angeben (6P)}
\begin{enumerate}[label=\alph*)]
    \item Tengo [Alter] años.
    \item Mi madre tiene [Alter] años.
    \item Mi hermano/hermana tiene [Alter] años.
\end{enumerate}
\textit{(Individuelle Antworten, Bewertung: korrekte tener-Form + Zahl + años)}

\textbf{Aufgabe 4: Familie beschreiben (6P)}
\begin{itemize}
    \item Mindestens 3 Sätze mit Familiengliedern und Alter
    \item Beispiel: ``Mi padre se llama Carlos. Tiene 45 años. Mi hermana se llama Ana. Tiene 10 años.''
\end{itemize}
\textit{Bewertung: Je 2P pro korrektem Satz (max. 6P)}

% ============================================================================
% 10. ANLAGEN
% ============================================================================
\section{Anlagen}

\begin{enumerate}
    \item Arbeitsblatt (AB-02b.pdf)
    \item Musterlösung (ML-02b.pdf)
    \item Lehrerhinweise (LH-02b.pdf)
    \item Lernpfad (LP-02b.html)
    \item Zahlenplakat 1-20
    \item Bingo-Karten (Kopiervorlage)
\end{enumerate}

% ============================================================================
% 11. DIDAKTIK-SCORE
% ============================================================================
\newpage
\section{Didaktik-Score}

\begin{scorebox}
\textbf{Bewertung der didaktischen Qualität nach 10 Dimensionen}\\
\textbf{Ziel-Score: $\geq$ 9.0 (Premium-Qualität)}

\vspace{1em}
\begin{tabular}{|c|l|c|c|p{4cm}|}
\hline
\textbf{\#} & \textbf{Dimension} & \textbf{Gew.} & \textbf{Score} & \textbf{Notiz} \\
\hline
1 & Differenzierung & 15\% & 9/10 & [B]/[S]/[E] qualitativ verschieden \\
\hline
2 & Taxonomie (AFB) & 10\% & 9/10 & 25/50/25 erfüllt \\
\hline
3 & Lernziel-Alignment & 10\% & 10/10 & RP MV explizit referenziert \\
\hline
4 & Aktivierung & 10\% & 10/10 & TPR, Bingo, Dialoge \\
\hline
5 & Scaffolding & 10\% & 9/10 & Gestufte Hilfen, Zahlenplakat \\
\hline
6 & Feedback-Qualität & 10\% & 9/10 & LP: elaboriertes Feedback \\
\hline
7 & Sprachliche Klarheit & 10\% & 10/10 & Kl. 7 angemessen \\
\hline
8 & Motivation & 10\% & 10/10 & Alltagsbezug (Geburtstag), Spiele \\
\hline
9 & Inklusion (UDL) & 10\% & 9/10 & Mehrere Zugänge (visuell, auditiv, kinästhetisch) \\
\hline
10 & Praktikabilität & 5\% & 9/10 & 90 Min. realistisch, Puffer vorhanden \\
\hline
\multicolumn{3}{|r|}{\textbf{GESAMT:}} & \textbf{9.3/10} & \\
\hline
\end{tabular}

\vspace{1em}
$\checkmark$ \textcolor{successgreen}{\textbf{FREIGABE} (Score $\geq$ 9.0)}
\end{scorebox}

\end{document}
